\section{Empirisk demonstrasjon}
\label{sec:empirical-nn}

Proposisjonane i teorien er ikkje berre logisk koherente; dei er empirisk falsifiserbare. Me testar kjernepredikasjonane mot eit korpus av europeiske stolar som spenner 744 år.

\subsection{Data}

Korpuset består av 2\,066 stolar (1280--2024) henta frå to museumssamlingar: Nasjonalmuseet, Oslo ($n = 63$) og Victoria and Albert Museum, London ($n = 2\,003$). Kvar post inkluderer katalogdimensjonar (høgde, breidde, djupn i centimeter), materialsamansetjing (fleirmateriell streng, t.d.\ ``eik, rotting, messing''), stilperiode (25 kanoniske kategoriar frå barokk til postmodernisme), nasjonalitet og dateringsintervall.

Katalogdataa er supplerte med 2\,204 tredimensjonale mesh-modellar med seks utrekna geometriske trekk: sfærisitet, fyllingsgrad, tregleiksratio, kompleksitet, konveks hylster-volum og mesh-vertekstal. Analysevinduet for tidsserietestar er 1500--2024, noko som gjev 1\,105 stolar med fullstendige dimensjonar og dateringsdata.

\subsection{Metodar}

Kvar proposisjon genererer ein testbar predikasjon. Me operasjonaliserer desse gjennom fem statistiske metodar:

\begin{enumerate}
\item \textbf{Gjensidig informasjon} (MI) samanliknar prediktiv styrke for stilperiode mot materiale på katalogdimensjonar (H, B, D, H/B). Stil er ein aggregert variabel som fangar mange samtidige trykk; materiale er eitt mekanisk trykk. Høgare MI for stil er difor delvis forventa \textit{a priori} og provar ikkje i seg sjølv at stil \textit{forårsakar} meir variasjon. Likevel er det konsistent med proposisjon 2.

\item \textbf{Variasjonskoeffisient} (CV) over dei seks mesh-trekka testar for eit avgrensningshierarki: dersom landskapet kanaliserer somme geometriske eigenskapar sterkare enn andre, bør CV-spreiinga vere stor (proposisjon 3).

\item \textbf{Silhouette-analyse} i firedimensjonalt mesh-rom testar om stilperiodar utgjer diskrete morfologiske klynger. Negative silhouette-skårar indikerer at stilkategoriar ikkje svarar til kompakte, velskilte klynger i formrommet; dette er konsistent med gradienthypotesen (proposisjon 3), men utelukkar ikkje alternative klyngjeløysingar.

\item \textbf{Kumulativt konveks hylster-volum} over suksessive 50-årsperiodar testar om det busette morforommet veks monotont (proposisjon 4). Monoton vekst er forventa for eitkvart akkumulerande utval; den informative samanlikninga er \textit{vekstraten} relativt til ein tilfeldig-trekking-nullmodell.

\item \textbf{Wasserstein-distanse} mellom konsekutive 50-årsperiodar testar om landskapet endrar seg signifikant over tid (proposisjon 4).
\end{enumerate}

Alle testar er kryssvaliderte mot 14 uavhengige delmengder: heile korpuset, kvart museum separat, sju hundreår-drop hold-outs, og to stil-drop hold-outs. Dette gjev 84 test-delmengd-kombinasjonar.

\subsection{Resultat}

Tabell~\ref{tab:evidence} oppsummerer dei ti kjernefunna.

\textbf{Morforommet er ikkje-uniformt busett} (proposisjon 1). Nærmaste-nabo variasjonskoeffisienten er 5,36; femten gonger Poisson-nullforventninga på 0,36 ($\mathrm{CI}_{95}$: [3,71; 5,94]).

\textbf{Stil predikerer meir enn materiale} (proposisjon 2). Gjensidig informasjon mellom stilperiode og form er 1,9--2,3 gonger større enn mellom materiale og form, over alle fire målte dimensjonar. Omfanget og konsistensen i gapet er kompatibelt med påstanden om at formvariasjon innanfor ein klasse er driven primært av krefter utover materiell avgrensing.

\textbf{Stilar utgjer ikkje diskrete klynger} (proposisjon 3). Silhouette-analyse i firedimensjonalt mesh-rom gjev ein skåre på $-0{,}338$ ($\mathrm{CI}_{95}$: [$-0{,}346$; $-0{,}329$]), med 21 av 25 stilkategoriar med negative individuelle skårar. Dette er konsistent med gradienthypotesen.

\textbf{Eit avgrensningshierarki finst} (proposisjon 3). Variasjonskoeffisienten spenner 128-falt over dei seks mesh-trekka ($\mathrm{CI}_{95}$: [50$\times$; 158$\times$]). Sfærisitet (CV $= 0{,}074$) er stramt avgrensa; volum (CV $= 9{,}39$) fluktuerer fritt.

\textbf{Morforommet ekspanderer kumulativt} (proposisjon 4). Det kumulative konvekse hylsteret ekspanderer med ein faktor på 553 i mesh-rom over ti 50-årsperiodar ($\mathrm{CI}_{95}$: [435$\times$; 752\,308$\times$]). Vekstraten overgår vesentleg ein tilfeldig-trekking-null.

\textbf{Landskapet endrar seg signifikant mellom periodar} (proposisjon 4). Gjennomsnittleg Wasserstein-1-distanse mellom konsekutive periodar er 14,4~cm (H), 8,2~cm (B) og 5,9~cm (D). Ingen av ti periodeovergangar viser nær-null-distanse.

\subsection{Robustheit og falsifisering}

Alle 84 test-delmengd-kombinasjonar produserer resultat konsistente med heile-korpus-funna. Desse er ikkje 84 uavhengige testar; delmengdene overlappar vesentleg. Konsistensen på tvers av medvite adversarielle hold-outs (museumsskifte, hundreårs-fjerning, stil-fjerning) gjev likevel evidens for at funna ikkje er artefakt av éi enkelt samling, periode eller stilkategori.

Sju direkte falsifiseringstestar målrettar teoriens postulat. Alle sju held med $p$-verdiar under $10^{-63}$ for dei sterkaste testane. Ingen postulat er falsifisert i datasettet. Fire tilleggstestar vart forsøkte og droppa for utilstrekkeleg effektstorleik; desse er dokumenterte transparent som negative eller inkonklusive resultat.
